\chapter{Hito 2: Instalación y Configuración de PostgreSQL}

\section{Objetivo}

El objetivo de este hito es preparar el motor de base de datos que servirá como núcleo del Data Warehouse macrofinanciero.

Al finalizar este hito el estudiante será capaz de:

\begin{itemize}
\item Instalar PostgreSQL.
\item Crear una base de datos para el proyecto.
\item Conectarse desde VS Code.
\item Configurar variables de entorno.
\item Conectarse desde Python.
\item Ejecutar consultas SQL básicas.
\end{itemize}

\section{Resultado esperado}

Al finalizar este hito deberá existir:

\begin{itemize}
\item Una instancia funcional de PostgreSQL.
\item Una base de datos llamada \texttt{macrofinanzas}.
\item Un usuario con permisos adecuados.
\item Un archivo \texttt{.env} para gestionar credenciales.
\item Un script Python capaz de conectarse a la base de datos.
\end{itemize}

\section{Paso 1: Verificar instalación de PostgreSQL}

Abrir PowerShell y ejecutar:

\begin{lstlisting}[language=bash]
psql --version
\end{lstlisting}

Salida esperada:

\begin{verbatim}
psql (PostgreSQL) 17.x
\end{verbatim}

Si el comando no funciona:

\begin{itemize}
\item Verificar que PostgreSQL está instalado.
\item Verificar que el directorio \texttt{bin} está en el PATH.
\end{itemize}

\section{Paso 2: Iniciar PostgreSQL}

Verificar que el servicio está activo.

En Windows:

\begin{enumerate}
\item Abrir Servicios.
\item Buscar PostgreSQL.
\item Verificar estado "Running".
\end{enumerate}

Alternativamente:

\begin{lstlisting}[language=bash]
net start postgresql-x64-17
\end{lstlisting}

\section{Paso 3: Conectarse al servidor}

Abrir terminal:

\begin{lstlisting}[language=bash]
psql -U postgres
\end{lstlisting}

Introducir la contraseña configurada durante la instalación.

Debería aparecer:

\begin{verbatim}
postgres=#
\end{verbatim}

\section{Paso 4: Crear la base de datos}

Ejecutar:

\begin{lstlisting}[language=SQL]
CREATE DATABASE macrofinanzas;
\end{lstlisting}

Verificar:

\begin{lstlisting}[language=SQL]
\l
\end{lstlisting}

La nueva base deberá aparecer en la lista.

Salir:

\begin{lstlisting}[language=SQL]
\q
\end{lstlisting}

\section{Paso 5: Conectarse a la nueva base}

\begin{lstlisting}[language=bash]
psql -U postgres -d macrofinanzas
\end{lstlisting}

Verificar:

\begin{lstlisting}[language=SQL]
SELECT current_database();
\end{lstlisting}

Resultado esperado:

\begin{verbatim}
macrofinanzas
\end{verbatim}

\section{Paso 6: Crear archivo .env}

En la raíz del proyecto crear:

\begin{verbatim}
.env
\end{verbatim}

Contenido:

\begin{lstlisting}
DB_HOST=localhost
DB_PORT=5432
DB_NAME=macrofinanzas
DB_USER=postgres
DB_PASSWORD=tu_password
\end{lstlisting}

\textbf{Importante:}

El archivo \texttt{.env} nunca debe enviarse a GitHub.

\section{Paso 7: Actualizar .env.example}

Crear una versión pública:

\begin{verbatim}
.env.example
\end{verbatim}

Contenido:

\begin{lstlisting}
DB_HOST=localhost
DB_PORT=5432
DB_NAME=macrofinanzas
DB_USER=postgres
DB_PASSWORD=your_password
\end{lstlisting}

Este archivo sí puede subirse al repositorio.

\section{Paso 8: Crear módulo de conexión}

Crear:

\begin{verbatim}
etl/db.py
\end{verbatim}

\begin{lstlisting}[language=Python]
from sqlalchemy import create_engine
from dotenv import load_dotenv
import os

load_dotenv()

HOST = os.getenv("DB_HOST")
PORT = os.getenv("DB_PORT")
DB = os.getenv("DB_NAME")
USER = os.getenv("DB_USER")
PASSWORD = os.getenv("DB_PASSWORD")

connection_string = (
f"postgresql+psycopg2://"
f"{USER}:{PASSWORD}@{HOST}:{PORT}/{DB}"
)

engine = create_engine(connection_string)

print("Conexion creada correctamente.")
\end{lstlisting}

\section{Paso 9: Probar la conexión}

Ejecutar:

\begin{lstlisting}[language=bash]
python etl/db.py
\end{lstlisting}

Salida esperada:

\begin{verbatim}
Conexión creada correctamente.
\end{verbatim}

\section{Paso 10: Realizar una consulta desde Python}

Modificar temporalmente:

\begin{lstlisting}[language=Python]
from sqlalchemy import text

with engine.connect() as conn:
resultado = conn.execute(
text("SELECT version();")
)

```
for fila in resultado:
    print(fila)
```

\end{lstlisting}

Ejecutar nuevamente:

\begin{lstlisting}[language=bash]
python etl/db.py
\end{lstlisting}

Resultado esperado:

\begin{verbatim}
PostgreSQL 17.x ...
\end{verbatim}

\section{Paso 11: Instalar extensiones de VS Code}

Instalar:

\begin{itemize}
\item PostgreSQL
\item SQLTools
\item SQLTools PostgreSQL Driver
\end{itemize}

Estas extensiones facilitarán la ejecución de consultas y exploración de objetos.

\section{Paso 12: Crear primer script SQL}

Crear:

\begin{verbatim}
sql/test_connection.sql
\end{verbatim}

Contenido:

\begin{lstlisting}[language=SQL]
SELECT current_database();

SELECT current_user;

SELECT version();
\end{lstlisting}

Ejecutarlo desde VS Code o psql.

\section{Checklist de validación}

Antes de continuar verificar:

\begin{itemize}
\item[$\Box$] PostgreSQL instalado.
\item[$\Box$] Servicio iniciado.
\item[$\Box$] Base de datos creada.
\item[$\Box$] Archivo .env configurado.
\item[$\Box$] Python conectado a PostgreSQL.
\item[$\Box$] Consulta SQL ejecutada correctamente.
\end{itemize}

\section{Entregable}

Al finalizar este hito el estudiante deberá disponer de:

\begin{itemize}
\item Base de datos \texttt{macrofinanzas}.
\item Módulo de conexión Python.
\item Variables de entorno configuradas.
\item VS Code conectado a PostgreSQL.
\item Primer script SQL ejecutado exitosamente.
\end{itemize}

\section{Próximo hito}

En el Hito 3 construiremos la primera capa de datos:

\begin{itemize}
\item Diseño de las fuentes.
\item Creación de datos simulados.
\item Organización de carpetas \texttt{raw} y \texttt{processed}.
\item Primera ingesta hacia PostgreSQL.
\end{itemize}
